\documentclass[10pt]{article}

\usepackage[utf8]{inputenc}

\usepackage[T1]{fontenc}

\usepackage{amsmath}

\usepackage{amsfonts}

\usepackage{amssymb}

\usepackage[version=4]{mhchem}

\usepackage{stmaryrd}

\usepackage{graphicx}

\usepackage[export]{adjustbox}

\graphicspath{ {./images/} }



\begin{document}

или описанной сферы. Тетраәдр двойствен сам себе, гексаэдр - октаэдру и додекаэдр - икосаэдру.



В пространстве $E^{4}$ существуют шесть $\Pi$. м., данные о к-рых приведевы в табл. 2 .



В пространстве $E^{n}, n>4$, существует три П. м.аналоги тетраэдра, октаэдра и куба; их символы ІІлефли $-\{3, \ldots, 3\},\{4,3, \ldots, 3\},\{3, \ldots, 3,4\}$.



Если под многоугольником понимать плоские замкнутые ломаные (хотя бы и самопересекающиеся), то можно указать еще 4 н е в ы п у к лы х (3 в е 3 д ч а-

\includegraphics[max width=\textwidth, center]{2023_07_08_ca2697b5cd1dd88fe416g-1}



Рис. 2.



т ы х) П. м. (т е л а П у а н с о). В этих многогранниках либо грани пересекают друг друга, либо грави самопересекающиеся многоугольники (рис. 2a-22). Данные о них приведены в табл. 3 .



та бл. 3. - П а вильны е (невы пуклы е)



\begin{center}

\includegraphics[max width=\textwidth]{2023_07_08_ca2697b5cd1dd88fe416g-1(1)}

\end{center}



Лит.: [1] Энциклопедия әлементарной математики, кн. 4Геометрия, М., 1963; [2] Л ю с т ер н и к Л. А., Выпуклые фигуры и многогранники, М., 1956; [3] Iї К л я р с к и й Д. О., Ч е в ц о в Н. Н., Я г л о м И. М., Избранные задачи и теоремы элементарной математики, ч. 3, M., $1954 ;$

x e t e r H. S. M., Regular polytopes, 3 ed., N.Y., 1973.



ПРАВИЛЬНЫЕ МНоГОУГОЛЬНИКИ - многоугольники, у к-рых равны все его углы и равны все его стороны. Подробнее см. Многоугольник.



ПРАВОУПОРЯДОЧЕННАЯ ГРУППА - групाа $G$, на множестве элементов которой задано отношение линейного порядка $\leqslant$ такое, что для всех $x, y, z$ из $G$ неравенство $x \leqslant y$ влечет за собой $x z \leqslant y z$. Множество $P=\{x \in G \mid x>e\}$ положительных элементов группы $G$ является ч и с т о й (то есть $\left.P \cap P^{-1}=\varnothing\right)$ л ине й но й (то есть $P \cup P^{-1} \cup\{e\}=G$ ) пол у гр у пп о й. Всякая чистая линейная подполугруппа $P$ произвольной группы определяет в ней правый порядок, а именно порядок $x<y \Longleftrightarrow y x^{-1} \in P$.



Группа $A(X)$ автоморфизмов линейно упорядоченного множества $X$ естественным образом может быть правоупорядочена. Всякая П. г. порядково изоморфна нек-рой подгруппе $A(X)$ для подходящего линейно упорядоченного множества (см. [1]). Архимедова П. г., то есть П. г., для к-рой верна аксиома Архимеда (см. $A p$ химедова группа), порядково изоморфна подгруппе аддитивной группы действительных чисел. В отличие от (двусторонне) линейно упорядоченных групп существуют некоммутативные П. г. без собственных выпуклых подгрупі. Класс П. г. замкнут относительно лексикографич. расширений. Система всех выпуклых подгрупп П. г. $G$ линейно упорядочена по включению и полна. Эта система разрешима тогда и только тогда, когда для всяких положительных әлементов $a, b \in G$ существует натуральное число $n$ такое, что $a^{n} b>a$. Если группа обладает разрешимой системой подгрупі $S(G)$, факторы к-рой не имеют кручения, то $G$ можно так правоупорядочить, что все подгруппы из $S(G)$ окажутся выпуклы- ми. В локально нильпотентной П. г. система выпуклых подгрупп разрешима.



Группа $G$ тогда ■ только тогда может быть правоупорядочена, когда для любого конечного набора



\[

\left\{x_{i} \neq e, 1 \leqslant i \leqslant n\right\}

\]



элементов из $G$ найдутся числа $\varepsilon_{i}= \pm 1,1 \leqslant i \leqslant n$, такие, что полугруппа, порожденная множеством $\left\{x_{1}^{\varepsilon_{1}}, \ldots, x_{n}^{\varepsilon_{n}}\right\}$, не содержит единицы группы $G$.



Всякий структурный (решеточный) порядок группы есть пересечение нек-рых ее правых порядков (см. Структурно упорлдоченнал групnа).



Лит:: [1] К о ко р вв А. И., К 0 п ы т о в B. М., Линейно упорядочевR. B., Rhem tu 11 a A., Orderable groups, N.Y.-Basel, 1977.,



B. $M$. Konыmos. ференциальное уравнение крыла самолета конечного размаха. При выводе П. у. делаются предположения, которые позволяют считать каждый элемент крыла находящимся в условиях обтекания его плоскопараллельным потоком. Это дает возможность связать геометрич. характеристики крыла с его аэродинамич. свойствами. Так, полученное П. у. имеет вид



\[

\frac{F\left(t_{0}\right)}{B\left(t_{0}\right)}+\frac{1}{\pi} \int_{-a}^{a} \frac{F^{\prime}(t)}{t-t_{0}} d t=f\left(t_{0}\right),-a<t_{0}<a

\]



где $F$ - искомая функция, $F^{\prime}(t)=d F(t) / d t, B$ и $f-$ заданные функции $B(t)=c b(t), f(t)=v \omega(t)$, а несобственный интеграл понимается в смысле главного значения по Коши. Значения входящих в эти равенства величин следующие: $2 a$ - размах крыла, к-рое предполагается симметричным относительно плоскости $\mathrm{Ozz}$, причем направление оси $O z$ совпадает с направлением потока воздуха на бесконечности; $b(x)$ обозначает хорду профиля, к-рый соответствует абсциссе $x ; F(x)-$ циркуляцию воздушного потока вокруг этого профиля; $c$ - нек-рую постоянную; $v-$ скорость воздушного потока на бесконечности; а $\omega$ - функцию, зависящую от изогнутости профиля и перекручивания крыла (см. [1]). Исходя из экспериментальных данных, полагают, что $F(-a)=F(a)$.



П. у. решается в замкнутом виде лишь при весьма жестких предположениях. В общем случае удается П. у. свести к интегральному уравнению Фредгольма (см. [3]).



П. у. наз. по имени Л. Прандтля (L. Prandtl).



М. Лит.: [1]. Го лу у б е в В. В., Лекции по теория крыла, че- - л., аәролинамика идеальных жидкостей, пер. с англ., М.-ЈI., 1939; [3] М у с х ел и ш в и л и Н. И., Сингулярные интегральные уравнения, 3 изд., М.,, 1968. Б. В. Хведелидзе.



ПРАНДТЛЯ ЧИСЛО - один из критериев подобия тепловых процессов в жидкостях и газах. П. ч. зависит только от термодинамич. состояния среды. П. ч.



\[

\operatorname{Pr}=v / a=\mu c_{p} / \lambda

\]



где $\nu=\mu / \rho-$ кинематич. коэффициент вязкости, $\mu-$ динамич. коэффициент теплопроводности, $\rho$ - плотность, $\lambda$ - коэффициевт теплопроводности, $a=\lambda / \rho c_{p}-$ коэффициевт температуропроводности, $c_{p}-$ удельная теплоемкость среды при постоянном давлении.



П. ч. связано с др. критериями подобия - Пекле



\includegraphics[max width=\textwidth, center]{2023_07_08_ca2697b5cd1dd88fe416g-1(2)}

П. ч. наз. по имени Л. Прандтля (L. Prandtl).



По материалам одноименной статьи из БС Э-з. ПРЕВОСХОДНОЕ КОЛЬЦО - коммутативное нётерово кольцо, удовлетворяющее трем приводимым ниже аксиомам. Известно, что геометрические кальца обладают рядом качественных свойств, не присущих





\end{document}